\chapter{Diretrizes para a comunicação de conselheiras e trabalhadoras do CRP~SP nas redes sociais}

A presença do Conselho Regional de Psicologia de São Paulo (CRP~SP) nas
redes sociais dá-se no cumprimento das atribuições da autarquia de
orientar e informar a categoria sobre temas associados à Psicologia, à
atuação profissional e ao funcionamento do Sistema Conselhos de
Psicologia. Objetiva complementar os canais proprietários da
instituição, buscando ampliar seus pontos de contato e comunicação com
psicólogas de todo o estado de São Paulo, assim como participar dos
debates, em sociedade e com demais entidades, sobre a Psicologia como
ciência e profissão.

As conselheiras em vigência de mandato representam o CRP~SP perante a
sociedade a qualquer tempo. Por isso, suas manifestações públicas ---
palestras, entrevistas a veículos de imprensa, comentários em redes
sociais etc. --- devem ser pautadas pelos princípios dispostos nesta
\textbf{Política}.

Fica definido que:

\begin{enumerate}

	\item
	só serão permitidos perfis e páginas associados ao CRP~SP que sejam
	criados e mantidos pela instituição, ou seja, tenham caráter de
	oficiais. A autarquia reserva-se o direito de não corroborar perfis e
	páginas que se autointitulem associados ao Conselho e não sejam
	reconhecidos como perfis oficiais, tomando as devidas medidas quando
	necessárias;
	\item
	as contas oficiais do CRP~SP serão gerenciadas por equipe certificada
	e autorizada pelo CRP~SP, composta por profissionais da comunicação,
	trabalhadoras e trabalhadores da autarquia e psicólogas e psicólogos
	conselheiras e conselheiros;
	\item
	como forma de garantia da veracidade, exatidão, confiabilidade,
	qualidade e legalidade das informações e mensagens que circulam pelos
	canais e perfis oficiais do CRP~SP, as interações, nestes espaços,
	estarão sujeitas à moderação e filtragem, podendo ser excluídas as
	mensagens e os perfis que

	\begin{enumerate}

		\item
		usem linguagem inapropriada, obscena, caluniosa, grosseira, abusiva,
		difamatória, ofensiva ou de qualquer outra forma censurável;
		\item
		façam apologia a práticas ilícitas;
		\item
		incitem o ódio, a violência, o racismo ou façam discriminação de
		qualquer ordem;
		\item
		contenham ameaças, assédio, injúria, calúnia ou difamação ou
		configurem qualquer outra forma de ilícito penal;
		\item
		divulguem conteúdos na forma de \emph{spam} ou ``correntes'';
		\item
		caracterizem intuito comercial ou publicitário;
		\item
		sejam ininteligíveis ou fora de contexto;
		\item
		contenham propagandas político-partidárias;
		\item
		contenham \emph{links} suspeitos ou representem ameaça à segurança
		da informação;
		\item
		façam uso de informações e imagens de pessoas e instituições de modo
		indevido;
		\item
		contenham dados pessoais da autora, do autor ou de terceiros;
		\item
		violem os direitos de imagem e de propriedade intelectual;
		\item
		sejam fraudulentos (perfis falsos e/ou \emph{bots}) ou promovam
		conteúdo inverídico (\emph{fake news}).
	\end{enumerate}
\end{enumerate}

A não observância a essas regras poderá acarretar ações imediatas do
CRP~SP, independentemente de justificativa, consulta ou aviso prévio.
Conforme o conteúdo, as mensagens poderão ser encaminhadas à autoridade
responsável.

\clearpage
\section{Orientações gerais}

\textbf{O ciberespaço é um espaço público}. Todas as opiniões,
comentários e interações emitidas/ocorridas em redes sociais, portanto,
devem ser entendidas como manifestações públicas, sujeitas à
responsabilização na forma da lei.

A participação do Sistema Conselhos de Psicologia na arena pública,
tanto em sua dimensão virtual quanto física, é estritamente
institucional. Por isso, recomenda-se o máximo de cuidado a conselheiras
e trabalhadoras do CRP~SP antes de emitir qualquer posicionamento,
comentário ou opinião pessoal sobre assuntos que digam respeito ao
Sistema Conselhos e/ou à prática da Psicologia. Além disso,
\textbf{perfis pessoais são expressamente vedados de falar em nome do
	CRP~SP}. Menções à autarquia, a conselheiras e trabalhadoras, desde que
não violem questões de direitos autorais, privacidade e segurança da
informação, e que sejam feitas com respeito e responsabilidade, são
autorizadas.

Caso seja necessário emitir um posicionamento e/ou resposta oficial
sobre qualquer assunto, a Unidade de Comunicação deve ser acionada por
um solicitante legítimo. É fundamental que o conteúdo enviado esteja
previamente validado por uma comissão permanente ou especial ou por uma
unidade administrativa, garantindo que o tema seja tratado de forma
\textbf{técnica} e alinhada aos parâmetros institucionais.

Toda e qualquer comunicação pessoal deve ser tratada como tal. Mensagens
de caráter pessoal não devem ser tornadas públicas. \textbf{Havendo
	necessidade de divulgar o conteúdo de alguma mensagem de terceiros,
	certifique-se de obter uma autorização prévia da autora}. Tendo em vista
a possibilidade de divulgação, mantenha, em sua correspondência
institucional, o mesmo respeito e urbanidade que caracterizam as
manifestações públicas do CRP~SP.

\textbf{No engajamento com postagens não relacionadas ao Sistema
	Conselhos e/ou à prática da Psicologia, conselheiras e trabalhadoras
	devem se portar de maneira respeitosa e condizente com os princípios
	definidos nas normativas do Sistema Conselhos e, mais amplamente, com o
	ordenamento jurídico brasileiro}. Ressaltamos a importância de não
disseminar nenhum conteúdo que não proceda de fonte identificável,
confiável e verificável.

\section{Dicas de comunicação \emph{on-line}}

A seguir, listamos uma série de dicas adaptadas do
\href{https://www.gov.br/economia/pt-br/centrais-de-conteudo/publicacoes/guias-e-manuais/ManualdeRedesSociais.pdf}{Manual
	de redes sociais} do Ministério da Economia, para que sejam usadas como
referência.

Os conteúdos postados são sempre de ordem pessoal, mas, a partir do
momento em que a usuária definir seu local de trabalho, invariavelmente
terão também um teor institucional/profissional. Isso significa que há
algumas boas práticas simples que devem ser seguidas quando a
conselheira ou trabalhadora postar alguma imagem ou texto de seu local
de trabalho ou com referência a ele:

\begin{itemize}
\item
as informações escritas na rede são de responsabilidade da
conselheira/trabalhadora, mas atingem a todas as entidades --- pessoas
e instituições --- presentes nos conteúdos de suas mensagens;
\item
você é uma pessoa pública. Sempre que postar algo nas redes sociais,
entenda que o conteúdo da sua mensagem será visto por colegas, chefes,
clientes, fornecedores, parceiras institucionais e de negócio, amigas
e familiares;
\item
suas seguidoras/amigas vão confundir o seu ``eu'' pessoal com o seu
``eu'' profissional. Você pode não ser a porta-voz oficial da
instituição na qual trabalha, mas, a partir do momento em que deixar
público seu perfil numa rede social, será vista pelas demais usuários
como alguém que fala em nome da instituição;
\item
evite postar qualquer coisa que possa gerar danos à instituição em que
atue;
\item
escrever na rede é o mesmo que escrever em pedra. Escrever não é o
mesmo que falar: suas palavras ficam na \emph{web} e são indexadas
quase que instantaneamente por outras redes. Assim, mesmo que apague
um \emph{post} do qual tenha se arrependido, ele provavelmente já terá
sido indexado pelo Google e por outros \emph{sites}, eternizando-se na
internet e pondo-se ao alcance de todos os usuários. Pense antes de
publicar. Se for para se arrepender, arrependa-se antes de escrever;
\item
proteja-se. Uma crise envolvendo \emph{posts} em redes sociais feitos
por uma trabalhadora e prejudicando o órgão nunca tem somente a
instituição como alvo único, e pode acabar causando danos de imagem
também a suas funcionárias.
\item
nunca deixe de ser você. Como qualquer cidadã, você é livre para
pensar e expressar o que desejar, da forma que preferir. Mas, como
qualquer pessoa pública, tem de entender que tudo o que expressar
provavelmente trará consequências, sejam positivas ou negativas;
\item
cuidado com questões sigilosas e informações privilegiadas. É vedada a
publicação de informação confidencial ou prejudicial à imagem do
Conselho, assim como não se deve antecipar resultados de trabalhos
ainda não divulgados oficialmente pela autarquia;
\item
tenha uma conduta ética nas redes sociais durante viagens de trabalho.
Conselheiras e/ou trabalhadoras que estiverem em viagem a serviço do
Conselho devem tomar cuidado com suas postagens em redes sociais. Há
diversos casos em que a imagem de órgãos públicos foi arranhada devido
a mau comportamento de servidores em viagens oficiais que é exposto
nas redes;
\item
lembre-se: você não está a turismo. Ainda que nas horas vagas você
aproveite uma ida pelo CRP~SP ao Rio de Janeiro, por exemplo, e
resolva pegar uma praia ou visitar um ponto turístico, mostrar isso
nas redes pode gerar um conflito de mensagens para seus seguidores e
principalmente para colegas de trabalho e amigos que sabem do real
propósito de sua missão. As viagens são custeadas com recursos
públicos e isso deve ser levado em máxima consideração;
\item
mostre a que veio. Postagens sobre o trabalho realizado, de forma
positiva, são salutares e até incentivadas. Divulgar o nosso trabalho
e o trabalho do CRP~SP é sempre positivo e colabora para a boa imagem
da instituição.
\end{itemize}